\clearpage
\VerdierPhysicalPageAnchor{0193}%
\begingroup\small\setlength{\parskip}{0pt}\setlength{\emergencystretch}{1.5em}
\setlength{\abovedisplayskip}{4pt plus 2pt minus 2pt}%
\setlength{\belowdisplayskip}{4pt plus 2pt minus 2pt}%

\VerdierSemanticAnchor{ch3d-2-2-6-proof-cont}
(2.1.1). The morphism $p_0$ is an elementary morphism of height $1$.
The complex $U_0^\bullet$ has the form:

\VerdierSemanticAnchor{ch3d-2-2-6-u0-complex}
\[
\begin{aligned}
U_0^\bullet={}&0\longrightarrow0\longrightarrow\cdots\longrightarrow0
\longrightarrow L_n\longrightarrow L_{n-1}\longrightarrow\cdots\\
&\longrightarrow L_1\longrightarrow Z_0\longrightarrow L
\longrightarrow0\longrightarrow0\longrightarrow\cdots.
\end{aligned}
\]
By property $(E_4)$, the object $Z_0$ belongs to $M'$. Consequently,
there is an epimorphism $u_0\colon L'_0\longrightarrow Z_0$ with
$L'_0\in M$. Set, as in (2.1.1),
$U_1^\bullet=U_0^\bullet(u_0,0)$ and
$p_1=p(u_0,0)\colon U_1^\bullet\longrightarrow U_0^\bullet$.

\VerdierSemanticAnchor{ch3d-2-2-6-iterated-elementary-morphisms}
Continuing in this way, induction gives a sequence of morphisms
\[
p_i\colon U_i^\bullet\longrightarrow U_{i-1}^\bullet,
\qquad 1\leq i\leq n-1,
\]
\[
p_0\colon U_0^\bullet\longrightarrow V^\bullet,
\]
where $p_i$ is an elementary morphism of height $-i+1$ and
$U_i^\bullet$ has the form:

\VerdierSemanticAnchor{ch3d-2-2-6-ui-complex}
\[
\begin{aligned}
U_i^\bullet={}&0\longrightarrow0\longrightarrow\cdots
\longrightarrow L_n\longrightarrow L_{n-1}\longrightarrow\cdots
\longrightarrow L_{i+1}\longrightarrow Z_i\\
&\longrightarrow L'_{i-1}\longrightarrow\cdots\longrightarrow L'_0
\longrightarrow L\longrightarrow0\longrightarrow\cdots ,
\end{aligned}
\]
where $L'_j\in M$, $0\leq j\leq i-1$. Composing these morphisms $p_i$
therefore gives an epimorphic quasi-isomorphism

\VerdierSemanticAnchor{ch3d-2-2-6-composite}
\[
p\colon U_{n-1}^\bullet\longrightarrow V^\bullet.
\]
Now observe that in the complex

\VerdierSemanticAnchor{ch3d-2-2-6-unminus1-complex}
\[
\begin{aligned}
U_{n-1}^\bullet={}&0\longrightarrow0\longrightarrow\cdots
\longrightarrow L_n\longrightarrow Z_{n-1}\longrightarrow L'_{n-2}
\longrightarrow\cdots\\
&\longrightarrow L'_0\longrightarrow L\longrightarrow0
\longrightarrow\cdots ,
\end{aligned}
\]
the object $Z_{n-1}$ appears, by property $(E_4)$, as an extension of two
objects of $M$ and hence belongs to $M$ by property $(E_5)$. Set

\VerdierSemanticAnchor{ch3d-2-2-6-l-complexes}
\[
L^\bullet=0\longrightarrow0\longrightarrow L_n\longrightarrow L_{n-1}
\longrightarrow\cdots\longrightarrow L_0\longrightarrow0
\longrightarrow0\longrightarrow\cdots ,
\]
\[
\begin{aligned}
{L'}^\bullet={}&0\longrightarrow0\longrightarrow L_n
\longrightarrow Z_{n-1}\longrightarrow L'_{n-2}
\longrightarrow\cdots\\
&\longrightarrow L'_0\longrightarrow0\longrightarrow0
\longrightarrow\cdots ,
\end{aligned}
\]
and let $m\colon {L'}^\bullet\longrightarrow L^\bullet$ be the morphism
induced by $p$. The only nonzero cohomology object of $L^\bullet$
(respectively, ${L'}^\bullet$) is $Y$ in degree $0$ (respectively, $L$ in
degree $0$), and the morphism induced by $m$ on these objects is $u$.
Consider the cone $c(m)$ of the morphism $m$:

\VerdierSemanticAnchor{ch3d-2-2-6-cone}
\[
\begin{aligned}
c(m)={}&\cdots\longrightarrow0\longrightarrow L_n
\longrightarrow L_n\oplus Z_{n-1}
\longrightarrow L_{n-1}\oplus L'_{n-2}\longrightarrow\cdots\\
&\longrightarrow L_1\oplus L'_0\longrightarrow L_0
\longrightarrow0\longrightarrow\cdots.
\end{aligned}
\]
\endgroup
